\section*{第 \ref{ch:g11:vision-and-brain} 章——视觉与大脑}

\begin{solution}{exo:g11:vision-and-brain:1}
来自左侧视野的光落在每个视网膜的右半部分；
左眼右半部分的纤维在视交叉处交叉，右眼的则留在同侧；
二者共同组成右侧视束，
在丘脑换元，
最终到达右侧\omterm{prop:g11:vision-and-brain:pathway}{初级视觉皮层}。
\end{solution}

\begin{solution}{exo:g11:vision-and-brain:2}
两条视神经在此相会；来自每侧视网膜鼻侧半的纤维交叉到对侧，
颞侧半的纤维则不交叉。
此后每条视束各携带视野的一半。
\end{solution}

\begin{solution}{exo:g11:vision-and-brain:3}
颜色区（所见世界为灰色）、
运动区（无法感知运动）、面孔区（无法识别面孔）。
\end{solution}

\begin{solution}{exo:g11:vision-and-brain:4}
可塑性：神经元连接能够根据使用情况
被增强、减弱、形成或修剪的能力。关键期：
生命早期的一个时期，在此期间，
经验对皮层的塑造作用决定性且不可逆。
\end{solution}

\begin{solution}{exo:g11:vision-and-brain:5}
它与视觉区神经元上的 5-羟色胺受体结合，
干扰信号的传递，
在没有相应光刺激的情况下产生幻觉。
\end{solution}

\begin{solution}{exo:g11:vision-and-brain:6}
视交叉之后，位于左侧的：左侧视束、
左侧丘脑中继或左侧\omterm{prop:g11:vision-and-brain:pathway}{初级视觉皮层}。
\end{solution}

\begin{solution}{exo:g11:vision-and-brain:7}
交叉的纤维来自两侧视网膜的鼻侧半，它们负责视野的外侧（颞侧）半，
因此患者丧失了左眼视野的左半
和右眼视野的右半，
表现为双侧隧道视——视野两侧均收窄。
\end{solution}

\begin{solution}{exo:g11:vision-and-brain:8}
正常：$12 + 20 + 36 + 20 = 88\%$。剥夺组：$2 + 3 + 5 + 30 = 40\%$，
其中主要的仅占 10\%。
\end{solution}

\begin{solution}{exo:g11:vision-and-brain:9}
遮盖疗法迫使皮层利用弱视眼的输入，
因此在关键期内这些连接得以保留并增强，而不是被修剪。
对成人而言，关键期已结束：
连接不再在这种尺度上重组，
因此弱视眼无法重新夺回其皮层区域。
\end{solution}

\begin{solution}{exo:g11:vision-and-brain:10}
她的\omterm{def:g11:the-eye:photoreceptors}{视锥细胞}及其\omterm{def:g11:the-eye:photoreceptors}{视蛋白}正常；
失去的是解读其比例的皮层阶段。
视网膜对彩色光反应的电生理记录在她身上应当正常，在色盲者身上则异常；
而她的缺陷是完全的且后天获得的，
色盲者的则是部分的且先天的。
\end{solution}

\begin{solution}{exo:g11:vision-and-brain:11}
多年的练习加强并增多了支配她手指的连接，
使其皮层区域扩大。
若停止练习，该区域会在数年间缩小，
正如杂耍研究所示。
\end{solution}

\begin{solution}{exo:g11:vision-and-brain:12}
起初她能看见光、颜色和移动的斑块，但无法识别形状、
面孔或深度：视网膜工作正常，但皮层在关键期内
从未建立起解读其信号的连接。她将逐渐、部分地学会使用视觉，
但永远达不到
自出生起就能看见的人那种熟练程度。
\end{solution}

\begin{solution}{exo:g11:vision-and-brain:13}
在中央凹处，每个\omterm{def:g11:the-eye:photoreceptors}{视锥细胞}都有自己的神经节\omterm{def:g10:cells-common-unit:cell}{细胞}，
因此视野中心每度发出的纤维远多于周边，
周边则是一百个\omterm{def:g11:the-eye:photoreceptors}{视杆细胞}共用一个神经节\omterm{def:g10:cells-common-unit:cell}{细胞}；皮层按纤维数分配面积，
因而按细节而非视野度数分配。
\end{solution}

\begin{solution}{exo:g11:vision-and-brain:14}
该区域的活动就是颜色感知，无论信号来自眼睛还是记忆；
感知是脑区所做的事，而非眼睛所输送的。幻觉
是由药物而不是由意图或光线引发的同一种活动，
即脑区的活动被药物启动。
\end{solution}

\begin{solution}{exo:g11:vision-and-brain:15}
眼睛把光转换为信号，看见则由大脑完成。
眼睛完好的患者，皮层损伤后可能丧失颜色、运动或面孔的感知；
眼睛完好的小猫，若皮层不允许其与一只眼睛建立连接，
则这只眼在功能上失明；
而某种药物既不作用于眼睛，也不改变皮层结构，
却能使人看见并不存在的事物。眼睛提供信号，大脑才能看见。
\end{solution}

\begin{solution}{pb:g11:vision-and-brain:1}
\textbf{1.} 左眼或其视神经，位于视交叉之前：
仅一只眼的纤维受到影响。

\textbf{2.} 右侧视束、丘脑或视觉皮层：两只眼对同一半视野
的缺损是共有的，只能位于两眼该半视野纤维
已汇合的地方，
即视交叉之后。

\textbf{3.} 视交叉本身：来自两眼视网膜鼻侧半的交叉纤维，
它们负责看见视野的颞侧半，
被切断。

\textbf{4.} 运动区位于初级皮层之外。视野完整与视敏度正常，
说明初级皮层及其之前的各级通路
均工作正常。

\textbf{5.} B、C 和 D 的视网膜完好；A 的视网膜则不确定。
可用检眼镜检查视网膜，并记录
其对闪光的电反应来核实。

\textbf{6.} 右侧皮层接收视野的左半部分。

\textbf{7.} 中央 $2^\circ$：\qty{1}{cm}；$2^\circ$ 条带位于
$40^\circ$ 处：\qty{0.1}{cm}。比值为 10。

\textbf{8.} 中央凹的\omterm{def:g11:the-eye:photoreceptors}{视锥细胞}排列紧密，每个\omterm{def:g11:the-eye:photoreceptors}{视锥细胞}都有各自的
神经节\omterm{def:g10:cells-common-unit:cell}{细胞}，因此中央视野发出的纤维远多于周边；
皮层为其分配的面积也按比例更大。

\textbf{9.} 左侧视野中央有一个小的盲点，对阅读影响严重；
此外，左侧视野周边还有一个较大的盲区，
较少被人注意。

\textbf{10.} 患者通过转动眼睛，把每个词移到视网膜中对应皮层
功能完好的区域上，借助损伤的边界；
速度慢，但可以做到。

\textbf{11.} $20 + 36 + 20 = 76\%$。

\textbf{12.} $2 + 3 + 5 + 30 = 40\%$，且仅 10\%
以左眼为主要或同等输入。

\textbf{13.} 其输入仍能到达皮层，但因缺少使用，
其连接被削弱并修剪；
右眼的连接则扩展进入了这一空间。

\textbf{14.} 该重组仅限于早期生命的关键期，
人类的对应现象是弱视，它只在儿童期发生，
此后便无法治疗。

\textbf{15.} 若两眼均不占优势，则不发生接管：
直方图大致保持对称，但神经元总体反应减少，
且小猫双眼视力都很差——
缺少竞争的剥夺会同时削弱双眼。

\textbf{16.} 大脑会抑制转向眼的图像，以避免复视；
该眼未被使用，其连接在关键期被修剪，
视力永久减弱。眼镜可矫正光学问题，但
该眼仍未被使用，因此修剪仍在继续。

\textbf{17.} 遮住好眼，迫使皮层使用弱眼的信号，
从而保留并加强其连接。

\textbf{18.} 关键期约在七岁结束；此后
皮层不再在这一尺度上重新分配区域，因此
弱眼丢失的连接无法重建。

\textbf{19.} 在成年个体中，连接已经建立且稳定；
另一只眼的接管是发育中的皮层在关键期内发生的现象，
而非成熟皮层的现象。

\textbf{20.} A：眼或视神经；C：视交叉；B：右侧视束、丘脑
或初级皮层；D：其后的运动区。皮层的可塑性
仅限于关键期，因此儿童的遮盖疗法
在四岁时见效，到十五岁时便会失效。
\end{solution}
